% GENERATED FILE. Edit canonical lessons, then run npm run generate:books.
% canonical-source-hash: fnv1a64:a334937a4da26d28
% canonical-lessons: GU-C403-dhorimarg, GU-C403-mail, GU-C403-kilomitar, GU-C403-phantavu, GU-C403-takravu

\chapter{On the Highway}
\label{ch:gu-on-the-highway}
\hlchaptermodality{403}

\begin{chapteropening}
\textbf{By the end of this chapter:} \emph{I can talk about distances on the highway.}

Complete the last lesson of chapter 403: I can talk about distances on the highway.
\end{chapteropening}

\section[\texorpdfstring{\gu{ધોરીમાર્ગ} (dhorīmārg)}{ધોરીમાર્ગ (dhorīmārg)}]{\gu{ધોરીમાર્ગ} (dhorīmārg) — a highway; a main road}

\label{lesson:GU-C403-dhorimarg}



\begin{quote}
Before the new one: say the Gujarati for a bus conductor, then the Gujarati for to drive a cart.
\end{quote}

\subsection*{You'll want to know: \gu{ધોરીમાર્ગ}}
\textbf{\gu{ધોરીમાર્ગ}} — \emph{dhorīmārg} — \textquotedblleft{}a highway; a main road\textquotedblright{}.

\begin{quote}
\gu{આ} \gu{ધોરીમાર્ગ} \gu{પર} \gu{ઘણી} \gu{ટ્રકો} \gu{ચાલે} \gu{છે}.
\end{quote}

\begin{quote}
\emph{ā dhorīmārg par ghaṇī ṭrako cāle che.}
\end{quote}

\begin{quote}
Many trucks run on this highway.
\end{quote}

A masculine noun.

\begin{cousinweb}
\textbf{\gu{ધોરીમાર્ગ}} (\emph{dhorīmārg}) is built from \gu{ધોરી}, \textquotedblleft{}main\textquotedblright{}, which CDIAL lists under Sanskrit \emph{dhaurēya}, \textquotedblleft{}a beast of burden; foremost\textquotedblright{} (CDIAL 6884), and Sanskrit \emph{mārga}, \textquotedblleft{}a road\textquotedblright{}: the road that carries the load.
\end{cousinweb}

\begin{grammarlens}[title={What you've built}]
One more word for getting around.
\end{grammarlens}

\subsection*{Your turn}
\begin{itemize}
\raggedright
  \item \emph{Say it:} \emph{dhorīmārg}
  \item \emph{Say it:} \emph{dhorīmārg}, once more
  \item \emph{Say it:} \emph{ā dhorīmārg par ghaṇī ṭrako cāle che.} — Many trucks run on this highway.
  \item \emph{Recall:} say the Gujarati for a bus conductor, then the Gujarati for to drive a cart, then say \emph{dhorīmārg} again
\end{itemize}

\begin{tcolorbox}[breakable,colback=teal!4,colframe=teal!35!black,title={Before you move on}]
What does \gu{ધોરીમાર્ગ} mean? (\textquotedblleft{}A highway; a main road\textquotedblright{}.) Say it once more.
\end{tcolorbox}
\section[\texorpdfstring{\gu{માઈલ} (māīl)}{માઈલ (māīl)}]{\gu{માઈલ} (māīl) — a mile}

\label{lesson:GU-C403-mail}



\begin{quote}
Before the new one: say the Gujarati for to drive a cart, then the Gujarati for a highway.
\end{quote}

\subsection*{You'll want to know: \gu{માઈલ}}
\textbf{\gu{માઈલ}} — \emph{māīl} — \textquotedblleft{}a mile\textquotedblright{}.

\begin{quote}
\gu{સ્ટેશન} \gu{અહીંથી} \gu{બે} \gu{માઈલ} \gu{દૂર} \gu{છે}.
\end{quote}

\begin{quote}
\emph{sṭeśan ahī̃thī be māīl dūr che.}
\end{quote}

\begin{quote}
The station is two miles from here.
\end{quote}

A masculine noun.

\begin{cousinweb}
\textbf{\gu{માઈલ}} (\emph{māīl}) is borrowed from English \emph{mile}.
\end{cousinweb}

\begin{grammarlens}[title={What you've built}]
One more word for getting around.
\end{grammarlens}

\subsection*{Your turn}
\begin{itemize}
\raggedright
  \item \emph{Say it:} \emph{māīl}
  \item \emph{Say it:} \emph{māīl}, once more
  \item \emph{Say it:} \emph{sṭeśan ahī̃thī be māīl dūr che.} — The station is two miles from here.
  \item \emph{Recall:} say the Gujarati for to drive a cart, then the Gujarati for a highway, then say \emph{māīl} again
\end{itemize}

\begin{tcolorbox}[breakable,colback=teal!4,colframe=teal!35!black,title={Before you move on}]
What does \gu{માઈલ} mean? (\textquotedblleft{}A mile\textquotedblright{}.) Say it once more.
\end{tcolorbox}
\section[\texorpdfstring{\gu{કિલોમીટર} (kilomīṭar)}{કિલોમીટર (kilomīṭar)}]{\gu{કિલોમીટર} (kilomīṭar) — a kilometre}

\label{lesson:GU-C403-kilomitar}



\begin{quote}
Before the new one: say the Gujarati for a highway, then the Gujarati for a mile.
\end{quote}

\subsection*{You'll want to know: \gu{કિલોમીટર}}
\textbf{\gu{કિલોમીટર}} — \emph{kilomīṭar} — \textquotedblleft{}a kilometre\textquotedblright{}.

\begin{quote}
\gu{રોજ} \gu{હું} \gu{પાંચ} \gu{કિલોમીટર} \gu{ચાલું} \gu{છું}.
\end{quote}

\begin{quote}
\emph{roj hũ pā̃c kilomīṭar cālũ chũ.}
\end{quote}

\begin{quote}
I walk five kilometres every day.
\end{quote}

A masculine noun.

\begin{cousinweb}
\textbf{\gu{કિલોમીટર}} (\emph{kilomīṭar}) is borrowed from English \emph{kilometre}.
\end{cousinweb}

\begin{grammarlens}[title={What you've built}]
One more word for getting around.
\end{grammarlens}

\subsection*{Your turn}
\begin{itemize}
\raggedright
  \item \emph{Say it:} \emph{kilomīṭar}
  \item \emph{Say it:} \emph{kilomīṭar}, once more
  \item \emph{Say it:} \emph{roj hũ pā̃c kilomīṭar cālũ chũ.} — I walk five kilometres every day.
  \item \emph{Recall:} say the Gujarati for a highway, then the Gujarati for a mile, then say \emph{kilomīṭar} again
\end{itemize}

\begin{tcolorbox}[breakable,colback=teal!4,colframe=teal!35!black,title={Before you move on}]
What does \gu{કિલોમીટર} mean? (\textquotedblleft{}A kilometre\textquotedblright{}.) Say it once more.
\end{tcolorbox}
\section[\texorpdfstring{\gu{ફંટાવું} (phanṭāvũ)}{ફંટાવું (phanṭāvũ)}]{\gu{ફંટાવું} (phanṭāvũ) — to branch off (of a road)}

\label{lesson:GU-C403-phantavu}



\begin{quote}
Before the new one: say the Gujarati for a mile, then the Gujarati for a kilometre.
\end{quote}

\subsection*{You'll want to know: \gu{ફંટાવું}}
\textbf{\gu{ફંટાવું}} — \emph{phanṭāvũ} — \textquotedblleft{}to branch off (of a road)\textquotedblright{}.

\begin{quote}
\gu{આગળ} \gu{જઈને} \gu{રસ્તો} \gu{બે} \gu{બાજુ} \gu{ફંટાય} \gu{છે}.
\end{quote}

\begin{quote}
\emph{āgaḷ jaīne rasto be bāju phanṭāy che.}
\end{quote}

\begin{quote}
Further on, the road branches off two ways.
\end{quote}

Past (it …): \gu{તે} \gu{રસ્તો} \gu{ફંટાયો}, \emph{te rasto phanṭāyo}.

\begin{cousinweb}
\textbf{\gu{ફંટાવું}} (\emph{phanṭāvũ}) goes back to \emph{*phāṇṭ}, \textquotedblleft{}to separate by beating or cutting\textquotedblright{}, a form CDIAL reconstructs; CDIAL lists the Gujarati verb, \textquotedblleft{}to be divided\textquotedblright{} (CDIAL 9071).
\end{cousinweb}

\begin{grammarlens}[title={What you've built}]
One more word for getting around.
\end{grammarlens}

\subsection*{Your turn}
\begin{itemize}
\raggedright
  \item \emph{Say it:} \emph{phanṭāvũ}
  \item \emph{Say it:} \emph{phanṭāvũ}, once more
  \item \emph{Say it:} \emph{āgaḷ jaīne rasto be bāju phanṭāy che.} — Further on, the road branches off two ways.
  \item \emph{Recall:} say the Gujarati for a mile, then the Gujarati for a kilometre, then say \emph{phanṭāvũ} again
\end{itemize}

\begin{tcolorbox}[breakable,colback=teal!4,colframe=teal!35!black,title={Before you move on}]
What does \gu{ફંટાવું} mean? (\textquotedblleft{}To branch off (of a road)\textquotedblright{}.) Say it once more.
\end{tcolorbox}
\section[\texorpdfstring{\gu{ટકરાવું} (ṭakrāvũ)}{ટકરાવું (ṭakrāvũ)}]{\gu{ટકરાવું} (ṭakrāvũ) — to crash (into); to collide head-on}

\label{lesson:GU-C403-takravu}



\begin{quote}
Before the new one: say the Gujarati for a kilometre, then the Gujarati for to branch off.
\end{quote}

\subsection*{You'll want to know: \gu{ટકરાવું}}
\textbf{\gu{ટકરાવું}} — \emph{ṭakrāvũ} — \textquotedblleft{}to crash (into); to collide head-on\textquotedblright{}.

\begin{quote}
\gu{ધુમ્મસમાં} \gu{બે} \gu{ટ્રક} \gu{ટકરાઈ}.
\end{quote}

\begin{quote}
\emph{dhummasmā̃ be ṭrak ṭakrāī.}
\end{quote}

\begin{quote}
Two trucks crashed into each other in the fog.
\end{quote}

Past (it …): \gu{તે} \gu{ટ્રક} \gu{ટકરાઈ}, \emph{te ṭrak ṭakrāī}.

\begin{cousinweb}
\textbf{\gu{ટકરાવું}} (\emph{ṭakrāvũ}) goes back to Sanskrit \emph{ṭakkarā}, \textquotedblleft{}a blow on the head\textquotedblright{} (CDIAL 5424); CDIAL lists the Gujarati verb, \textquotedblleft{}to collide\textquotedblright{}.
\end{cousinweb}

\begin{grammarlens}[title={What you've built}]
One more word for getting around.
\end{grammarlens}

\subsection*{Your turn}
\begin{itemize}
\raggedright
  \item \emph{Say it:} \emph{ṭakrāvũ}
  \item \emph{Say it:} \emph{ṭakrāvũ}, once more
  \item \emph{Say it:} \emph{dhummasmā̃ be ṭrak ṭakrāī.} — Two trucks crashed into each other in the fog.
  \item \emph{Recall:} say the Gujarati for a kilometre, then the Gujarati for to branch off, then say \emph{ṭakrāvũ} again
\end{itemize}

\begin{tcolorbox}[breakable,colback=teal!4,colframe=teal!35!black,title={Before you move on}]
What does \gu{ટકરાવું} mean? (\textquotedblleft{}To crash (into); to collide head-on\textquotedblright{}.) Say it once more.
\end{tcolorbox}
